% Teil II: der durchgerechnete Weg.

\subsection*{Gegeben und gesucht}
Gegeben: $m_1=\qty{10}{kg}$, $m_2=\qty{5}{kg}$, $R=\qty{20}{cm}=\qty{0.20}{m}$,
$\omega_1=60\,\mathrm{min}^{-1}=\qty{6.28}{rad/s}$,
$\omega_2=-150\,\mathrm{min}^{-1}=-\qty{15.7}{rad/s}$.
Gesucht: (a) gemeinsame Winkelgeschwindigkeit $\omega$; (b) Reibungsenergie.

\subsection*{Rechnung}
\textbf{Winkelgeschwindigkeiten in SI.}
$60\,\mathrm{min}^{-1}=1\,\mathrm s^{-1}\;\Rightarrow\;\omega_1=2\pi\cdot
1=\qty{6.283}{rad/s}$. $150\,\mathrm{min}^{-1}=2.5\,\mathrm s^{-1}\;\Rightarrow\;
\omega_2=-2\pi\cdot 2.5=-\qty{15.71}{rad/s}$.

\textbf{Trägheitsmomente:}
\begin{equation}
  I_1=\frac12 m_1R^2=\frac12\cdot10\cdot0.20^2=\qty{0.20}{kg\,m^2},
  \qquad
  I_2=\frac12\cdot5\cdot0.20^2=\qty{0.10}{kg\,m^2}.
\end{equation}

\textbf{a) Gemeinsame Winkelgeschwindigkeit} (Drehimpulserhaltung):
\begin{equation}
  \omega
  = \frac{I_1\omega_1+I_2\omega_2}{I_1+I_2}
  = \frac{0.20\cdot6.283 + 0.10\cdot(-15.71)}{0.30}\ \si{rad/s}
  = \frac{1.257-1.571}{0.30}
  = \qty{-1.05}{rad/s}.
\end{equation}
Das entspricht $\omega/(2\pi)=-\qty{0.167}{s^{-1}}$, also $-10\,\mathrm{min}^{-1}$:
die Scheiben drehen gemeinsam mit $\qty{1.05}{rad/s}$ in die Richtung der
oberen (schnelleren) Scheibe.

\textbf{b) Reibungsenergie:}
\begin{equation}
  \Delta E
  = \frac12 I_1\omega_1^2+\frac12 I_2\omega_2^2 - \frac12(I_1+I_2)\omega^2
\end{equation}
\begin{equation}
  = \frac12\cdot0.20\cdot6.283^2
  + \frac12\cdot0.10\cdot(-15.71)^2
  - \frac12\cdot0.30\cdot(-1.05)^2\ \si{J}
  = (3.95 + 12.34 - 0.16)\ \si{J}
  \approx \qty{16.1}{J}.
\end{equation}

\begin{ergebnis}
(a) $\omega=\dfrac{I_1\omega_1+I_2\omega_2}{I_1+I_2}
=\qty{-1.05}{rad/s}$ (Drehsinn der oberen Scheibe). (b)
$\Delta E_\mathrm{Reibung}\approx\qty{16.1}{J}$.
\end{ergebnis}

\begin{kontrolle}
\textbf{Einheiten:} $\omega$: $kg\,m^2\cdot rad/s / kg\,m^2 = rad/s$. \checkmark
$\Delta E$: $kg\,m^2/s^2\cdot(rad/s)^2\,s^2/(rad^2)=J$. \checkmark

\textbf{Drehimpulserhaltung gegen Energie:} $E_\mathrm{kin}$ muss abnehmen
(Reibung) - $16.1\,\si{J}>0$. \checkmark

\textbf{Grenzfall:} stünden beide Scheiben gleich schnell ($\omega_2=\omega_1$),
kehrte sich nichts um und $\Delta E=0$; hier kommt die Reibung von der
unterschiedlichen Drehrichtung. Plausibel.
\end{kontrolle}

\begin{merksatz}
Wenn zwei rotierende Körper reibend zu einer gemeinsamen Winkelgeschwindigkeit
zusammenkommen, ist der \emph{Drehimpuls} (nicht die Energie) erhalten.
Man berechnet zuerst $I=\frac12 mR^2$, dann $L$, und zieht die
Reibungsenergie aus der Energiebilanz.
\end{merksatz}